\chapter{HP--Li Positivity and the Rankin--Selberg Kernel Obstruction}
\label{v4-arith:w5-a:chapter}

The HP--Li programme reduces the Riemann Hypothesis to the
positivity of the Weil distribution \(W(f*\widetilde{f})\geq 0\) on the
Paley--Wiener class \(\mathrm{PW}(\mathbb{R})\)
(\Cref{v4-arith:w5-a:cor:weil-positivity-RH}). The
fourth input to the sufficient RH path of
Chapter~\ref{v4-arith:rh-reformulation}, HP--Li positivity, is the
full-strength Weil positive trace inequality.

A coefficientwise Mellin lift
\(\Upsilon_{F}\colon\mathrm{PW}^{\mathrm{hol}}_{F}\to q\mathbb C[[q]]\)
constructed from a holomorphic cuspidal Hecke eigenform \(F\) of weight
\(k\geq 12\) is a formal probe of the tempered cuspidal core. It is not,
for a general Paley--Wiener multiplier \(f(\log n)\), a modular form.
The Rankin--Selberg kernel attached to a fixed \(F\) is the automorphic
kernel of \(L(s,F\otimes\overline F)\), whereas the finite term in
Weil's formula is the von Mangoldt kernel of \(\zeta\). Thus the
prime-side Weil contribution is not made positive by this lift.

The archimedean contribution \(W_{\infty}(f*\widetilde{f})\) is the
pairing of \(|\widehat{f}|^{2}\) against the distribution
\(\mathcal{T}_{\infty}(t)=\frac{1}{2\pi}\mathrm{Re}\,\psi(\tfrac{1}{4}+\tfrac{it}{2})+\frac{\log\pi}{2\pi}\),
where \(\psi=\Gamma'/\Gamma=\partial_{s}\log\Gamma\) is the digamma function.
The Mellin--Petersson isometry pairs \(|\widehat{f}|^{2}\) against the
Petersson archimedean density \(|\Gamma(\tfrac{1}{4}+\tfrac{it}{2})|^{2}\);
the Weil archimedean term \(W_{\infty}\) pairs against the
logarithmic-derivative density \(\mathrm{Re}\,\psi(\tfrac{1}{4}+\tfrac{it}{2})\).
These two densities are not proportional, so \(W_{\infty}\) lies outside
the isometry.
Conjecture A-TP (\Cref{v4-arith:w5-a:conj:archimedean-tate-positivity})
is the full Weil inequality with the archimedean term isolated on one
side. It is strictly equivalent to Weil positivity and to RH
(\Cref{v4-arith:w5-a:thm:ATP-equals-weil-positivity}).

\section{Weil's Explicit Formula (Recall)}
\label{v4-arith:w5-a:weil-explicit}

\subsection{Paley--Wiener Test Class}
\label{v4-arith:w5-a:weil-explicit:PW}

\begin{definition}[admissible Paley--Wiener class]
\label{v4-arith:w5-a:def:PW}
\ClaimStatusDefinitional. Let \(\mathrm{PW}(\mathbb{R})\) denote the
space of entire functions \(g(s)\) of exponential type such that
\(g(s)=\widehat{f}(s)\) for some even Schwartz-class function
\(f\in\mathcal{S}(\mathbb{R})\) of compact support, with the Fourier
normalisation \(\widehat{f}(s)=\int_{\mathbb{R}}f(u)e^{isu}\,du\). For
\(f\in\mathrm{PW}(\mathbb{R})\), write \(\widetilde{f}(u):=\overline{f(-u)}\).
The convolution
\((f*\widetilde{f})(u)=\int_{\mathbb{R}}f(v)\widetilde{f}(u-v)\,dv\) has
\[
\widehat{f*\widetilde{f}}(s)=
\widehat f(s)\,\overline{\widehat f(\overline s)}.
\]
On the real axis this is \(|\widehat f(t)|^2\); off the real axis it is
not an absolute square. For \(f\) real even, \(\widetilde{f}=f\).
\end{definition}

\begin{theorem}[Weil 1952 explicit formula]
\label{v4-arith:w5-a:thm:weil-explicit}
\ClaimStatusProvedElsewhere. Let \(\psi=\Gamma'/\Gamma\) be the
digamma function. For \(f\in\mathrm{PW}(\mathbb{R})\) with
Fourier transform \(\widehat{f}\), write
\(\gamma_{\rho}:=-i(\rho-\tfrac12)\in\mathbb C\) for each non-trivial
zero \(\rho\) of \(\zeta\). The Weil distribution satisfies
\begin{equation}
\label{v4-arith:w5-a:eq:weil-explicit}
\sum_{\rho}\widehat{f}(\gamma_{\rho})
=\widehat{f}\bigl(\tfrac{1}{2i}\bigr)+\widehat{f}\bigl(-\tfrac{1}{2i}\bigr)
-W_{\infty}(f)-W_{\mathrm{fin}}(f),
\end{equation}
with archimedean and finite-place Weil terms
\begin{align}
\label{v4-arith:w5-a:eq:weil-archimedean}
W_{\infty}(f)&:=\int_{-\infty}^{\infty}\widehat{f}(t)\cdot\frac{1}{2\pi}\Bigl[\mathrm{Re}\,\psi\bigl(\tfrac{1}{4}+\tfrac{it}{2}\bigr)+\log\pi\Bigr]\,dt,\\
\label{v4-arith:w5-a:eq:weil-finite}
W_{\mathrm{fin}}(f)&:=2\sum_{p}\sum_{n\geq 1}(\log p)\,p^{-n/2}\,f(n\log p).
\end{align}
The two polar terms \(\widehat{f}(\pm 1/(2i))=\widehat{f}(\mp i/2)\)
arise from the simple poles of the completed zeta \(\Lambda_{\zeta}\) at
\(s=1,0\) (the pole of \(\zeta\) at \(s=1\) and of \(\Gamma(s/2)\) at
\(s=0\)); the factor \(s(s-1)/2\) in \(\xi_{\mathbb{R}}=\tfrac12 s(s-1)\Lambda_{\zeta}\)
cancels them, so \(\xi_{\mathbb{R}}\) itself is entire.
\end{theorem}

\begin{proof}
Weil 1952 Selecta (\emph{Sur les formules explicites de la th\'eorie
des nombres premiers}, reprinted in \emph{Oeuvres Scientifiques} Vol.~II)
computes the distribution by residue calculus applied to
\(-\frac{1}{2\pi i}\oint\widehat{f}\,d\log\xi_{\mathbb{R}}\) on a
symmetric contour; the trivial-zero gamma factors produce
\(W_{\infty}\) and the Euler product yields \(W_{\mathrm{fin}}\).
Modern references: Iwaniec--Kowalski \emph{Analytic Number Theory}
(AMS 2004) Thm.~5.12; Patterson \emph{An Introduction to the Theory of
the Riemann Zeta-Function} (Cambridge 1988) Thm.~3.1.
\end{proof}

\begin{corollary}[Weil positivity as RH]
\label{v4-arith:w5-a:cor:weil-positivity-RH}
\ClaimStatusProvedElsewhere. RH is equivalent to
\begin{equation}
\label{v4-arith:w5-a:eq:weil-positivity}
W(f*\widetilde{f})\;\geq\;0\qquad\text{for all }f\in\mathrm{PW}(\mathbb{R}),
\end{equation}
where \(W(g):=\widehat{g}(\tfrac{1}{2i})+\widehat{g}(-\tfrac{1}{2i})-W_{\infty}(g)-W_{\mathrm{fin}}(g)\)
denotes the right-hand side of \eqref{v4-arith:w5-a:eq:weil-explicit}.
Under RH, \(\gamma_{\rho}\in\mathbb R\) and
\(W(f*\widetilde{f})=\sum_{\rho}|\widehat{f}(\gamma_{\rho})|^{2}\geq 0\);
the converse is Weil's positivity equivalence.
\end{corollary}

\begin{proof}
Weil 1952 \S IV. The key identity is
\((f*\widetilde{f})\widehat{}(t)=|\widehat{f}(t)|^{2}\geq 0\); on the
critical line \(\mathrm{Re}\,s=1/2\) the Weil distribution becomes a
sum of non-negative values, and the off-critical locus contributes a
negative term by Hadamard product considerations.
\end{proof}

\subsection{Li Coefficients as Weil Distribution Values}
\label{v4-arith:w5-a:weil-explicit:Li}

\begin{theorem}[Li 1997, J.~Number Theory 65]
\label{v4-arith:w5-a:thm:li-criterion}
\ClaimStatusProvedElsewhere. RH is equivalent to \(\lambda_{n}\geq 0\)
for every integer \(n\geq 1\), where
\begin{equation}
\label{v4-arith:w5-a:eq:li}
\lambda_{n}\;=\;\sum_{\rho}\biggl(1-\Bigl(1-\frac{1}{\rho}\Bigr)^{n}\biggr),
\end{equation}
summed over non-trivial zeros with multiplicity.
\end{theorem}

\begin{proposition}[Li coefficients and the Paley--Wiener boundary]
\label{v4-arith:w5-a:prop:li-via-weil}
\ClaimStatusProvedHere. The Li coefficients are values of the
zero divisor against the rational functions \(1-(1-1/s)^n\). This is a
Weil-type distributional statement, not a compact-support Paley--Wiener
interpolation theorem. The existence of
\(\phi_n\in\mathrm{PW}(\mathbb R)\) with
prescribed values
\(\widehat{\phi_n}(-i(\rho-\tfrac12))=1-(1-1/\rho)^n\) at all
non-trivial zeros is an additional interpolation theorem, not a
consequence of Li's criterion.
\end{proposition}

\begin{proof}
Li's criterion and the Bombieri--Lagarias form give the displayed
rational test on the zero set. Paley--Wiener interpolation at the zero
divisor would require an exponential-type interpolation theorem with
growth constants matched to the Riemann--von Mangoldt density. Li's
criterion supplies the rational values on the divisor; it does not
produce a compactly supported Paley--Wiener representative.
\end{proof}

\section{Petersson on the Tempered Cuspidal Core (Recall)}
\label{v4-arith:w5-a:petersson}

\begin{theorem}[Petersson positivity and native transfer]
\label{v4-arith:w5-a:thm:petersson-positive}
The squared $L^2$ norm of a square-integrable automorphic
function is nonnegative and vanishes only on the zero $L^2$
class. If a native core admits an injective realization $I$
in this Hilbert space, its pullback
\begin{equation}
 \langle v,w\rangle^{\mathrm{Pet}}
       =\int (Iv)(g)\overline{(Iw)(g)}\,dg
 \label{v4-arith:w5-a:eq:petersson}
\end{equation}
is positive definite. Identification with a native Koszul form
requires the pairing-map hypotheses of
\Cref{v4-arith:kp:hyp:native-pairing-maps} in addition.
\end{theorem}
\begin{proof}
The first assertion is the defining norm property of $L^2$.
Injectivity of $I$ transfers it to the native core. The bilinear
Koszul comparison is then
$B(v,cw)=B_q(fv,J_qfw)=h_q(fv,fw)$, provided the displayed
native maps exist. A positive local Gibbs state does not itself
construct either $I$ or those maps.
\end{proof}

\begin{remark}[the exact local normalization]
\label{v4-arith:w5-a:rem:restricted-tensor}
At $\beta=1$, the GNS creation vectors have norm one, while
the division vectors have squared norms $m^{-1}$. The radial
shell model maps its basis of squared norm $m^{-1}$ to
$m^{-1/2}[\mu_m]$ and its dual basis to $[\mu_m^*]$;
\Cref{v4-arith:kp:radial-pairing-realization} proves the pairing
and conjugation identities. The full local GNS carrier is
Hilbert--Schmidt, not a zero-mode lattice times an independent
orthogonal phase coordinate. Tensoring the local shift states
also does not produce the native adelic phase comparison.
\end{remark}

\section{The Mellin Lift and Rankin--Selberg Kernel Check}
\label{v4-arith:w5-a:mellin-lift}

The prime-side Weil contribution \(W_{\mathrm{fin}}(f*\widetilde{f})\) of
\eqref{v4-arith:w5-a:eq:weil-finite} is a sum over prime powers of a
Paley--Wiener weight. A Petersson proof of its positivity would require
a linear map from \(\mathrm{PW}(\mathbb{R})\) to the tempered cuspidal
core whose Petersson norm computes the von Mangoldt prime-power sum.
The coefficientwise lift \(\Upsilon_F\), defined below, does not supply
such a map without an additional automorphic multiplier theorem.

\subsection{Definition of \(\Upsilon\)}
\label{v4-arith:w5-a:mellin-lift:def}

\begin{definition}[Mellin lift]
\label{v4-arith:w5-a:def:upsilon}
\ClaimStatusDefinitional. Fix a normalised holomorphic cuspidal
Hecke eigenform \(F\in S_{k}(\mathrm{SL}_{2}(\mathbb{Z}))\) of weight
\(k\geq 12\), with Fourier expansion
\(F(z)=\sum_{n\geq 1}a_{F}(n)\,e^{2\pi inz}\) and Hecke normalisation
\(a_{F}(1)=1\). For \(f\in\mathrm{PW}(\mathbb{R})\), define
\begin{equation}
\label{v4-arith:w5-a:eq:upsilon}
\Upsilon_{F}(f)(z)\;:=\;\sum_{n\geq 1}f(\log n)\cdot a_{F}(n)\cdot n^{(k-1)/2}\,e^{2\pi inz},
\end{equation}
a modular-type formal series on the upper half-plane.
\end{definition}

By Deligne's bound for holomorphic cuspidal \(F\) (Deligne 1974
Publ.~IHES~43), the normalised Hecke eigenvalues satisfy
\(|a_{F}(p)p^{-(k-1)/2}|\le 2\). This is a growth estimate for the
coefficient sequence; it is not an automorphy theorem for an arbitrary
coefficientwise multiplier \(f(\log n)\).

\begin{lemma}[coefficientwise lift is only a formal automorphic probe]
\label{v4-arith:w5-a:lem:upsilon-cuspidal}
\ClaimStatusProvedHere. Let
\(\mathrm{PW}^{\mathrm{hol}}_{F}\subset\mathrm{PW}(\mathbb{R})\) denote
the subspace of \(f\in\mathrm{PW}(\mathbb{R})\) satisfying:
\begin{itemize}
\item[(H1)] \(f\) is real-valued and even: \(\widetilde{f}=f\);
\item[(H2)] the Mellin transform
\(\widetilde{f}_{M}(s):=\int_{0}^{\infty}f(\log u)\,u^{s-1}\,du\) is
holomorphic and of polynomial decay on the strip
\(|\mathrm{Re}\,s-1/2|\leq(k-1)/2-\varepsilon\) for some \(\varepsilon>0\);
\item[(H3)] \(f\)'s support avoids the Tate-polar contribution in the
sense: \(\widetilde{f}_{M}(s)\) has a double zero at \(s=0\) and \(s=1\).
\end{itemize}
Then \(\Upsilon_F(f)\) is a convergent holomorphic \(q\)-series on the
upper half-plane under the stated decay hypotheses, with coefficients
bounded by Deligne's estimate. It is a modular form of weight \(k\) only
under the additional condition that the multiplier
\(a_F(n)\mapsto f(\log n)a_F(n)\) is induced by an endomorphism of the
automorphic representation generated by \(F\). Conditions (H1)--(H3) do
not imply this condition.
\end{lemma}

\begin{proof}
Under (H2) and Deligne's bound \(|a_{F}(n)|\leq\sigma_{0}(n)n^{(k-1)/2}\)
(Deligne 1974), the series
\eqref{v4-arith:w5-a:eq:upsilon} converges absolutely on the upper
half-plane to a holomorphic function with no constant term. Modularity
is not a coefficientwise condition. The functional equation of \(L(s,F)\)
controls the original eigenform \(F\); it does not imply that the
Hadamard product of its Fourier coefficients with the sequence
\(f(\log n)\) transforms under \(S=\begin{pmatrix}0&-1\\1&0\end{pmatrix}\).
Thus the lift is a formal probe until an automorphic multiplier theorem
is supplied.
\end{proof}

\begin{remark}[domain of the restricted class]
\label{v4-arith:w5-a:rem:domain-PW-hol}
\(\mathrm{PW}^{\mathrm{hol}}_{F}\) is a restricted test class inside
\(\mathrm{PW}(\mathbb{R})\). No density assertion for a fixed \(F\) is
used. Condition (H2) restricts the spectral width (the Paley--Wiener
support must sit inside the strip where the weight-\(k\) Mellin
transform is holomorphic), and (H3) removes the Tate polar locus.
(H2) is a \emph{genuine} restriction whose complement is not reached
by a fixed holomorphic \(F\)-lift
(\Cref{v4-arith:w5-a:sec:obstruction}). (H3) is removable
by absorbing the two Tate poles into the archimedean term via the
entire form \(\xi_{\mathbb{R}}(s)=\frac{1}{2}s(s-1)\Lambda_{\zeta}(s)\);
see \Cref{v4-arith:w5-a:rem:H3-removable}.
\end{remark}

\subsection{Rankin--Selberg Integral and the Kernel Mismatch}
\label{v4-arith:w5-a:mellin-lift:RS}

The Petersson norm of a genuine cusp form admits a Rankin--Selberg
expansion. That expansion is not the finite-place Weil contribution for
\(\zeta\).

\begin{theorem}[Rankin--Selberg kernel does not compute the zeta prime term]
\label{v4-arith:w5-a:thm:RS-integral}
\ClaimStatusProvedHere. Suppose that \(\Upsilon_F(f)\) is in fact a
weight-\(k\) cusp form. Its Petersson norm is governed by the
Rankin--Selberg convolution of the automorphic representation generated
by \(F\). In any such identity the spectral kernel is built from
\begin{equation}
\label{v4-arith:w5-a:eq:KF-kernel}
K_{F}(s)\;:=\;\frac{\zeta(2s)}{\zeta(s)\xi_{\mathbb{R}}(s)\xi_{\mathbb{R}}(1-s)}\cdot L(s,F\otimes\overline{F})
\end{equation}
up to normalising gamma and conductor factors. The finite-place term in
Weil's explicit formula is instead the von Mangoldt distribution
\[
W_{\mathrm{fin}}(\phi)=2\sum_{p}\sum_{n\geq1}(\log p)p^{-n/2}\phi(n\log p).
\]
No Rankin--Selberg identity identifies \(W_{\mathrm{fin}}(f*\widetilde f)\)
with a positive Petersson norm attached to a fixed \(F\).
\end{theorem}

\begin{proof}
Rankin's unfolding computes Dirichlet series with coefficients
\(|a_F(n)|^2\), hence \(L(s,F\otimes\overline F)\) after the usual
normalisation. Weil's finite distribution has coefficients
\(\Lambda(n)n^{-1/2}\) and is the logarithmic derivative of \(\zeta\).
These coefficient systems are different. Positivity of a Petersson norm
therefore cannot be transported to \(W_{\mathrm{fin}}\) without an
additional trace formula identifying the two kernels.
\end{proof}

\begin{remark}[convention absorbing the Tate gamma]
\label{v4-arith:w5-a:rem:gamma-absorbed}
The kernel \(K_{F}(s)\) in \eqref{v4-arith:w5-a:eq:KF-kernel} is
written with archimedean Tate gamma factors on both top and bottom to
exhibit the Hecke functional equation symmetrically. The
\(\xi_{\mathbb{R}}(s)\xi_{\mathbb{R}}(1-s)\) in the denominator is
part of the automorphic normalisation of the Rankin--Selberg kernel.
It is not the statement that the zeta prime distribution
\(W_{\mathrm{fin}}\) has been absorbed into the Petersson form.
\end{remark}

\section{The Prime-Side Obstruction}
\label{v4-arith:w5-a:prime-closure}

The comparison with Rankin--Selberg exposes a kernel mismatch.

\begin{theorem}[no prime-side positivity from a fixed \(F\)-lift]
\label{v4-arith:w5-a:thm:prime-side}
\ClaimStatusProvedHere. The identity
\[
W_{\mathrm{fin}}(f*\widetilde f)
=2\bigl\langle\Upsilon_F(f),\Upsilon_F(f)\bigr\rangle^{\mathrm{Pet}}_{F\text{-normalised}}
\]
does not follow from Rankin--Selberg unfolding for a fixed cusp form
\(F\). Consequently fixed-\(F\) Rankin--Selberg positivity does not imply
\(W_{\mathrm{fin}}(f*\widetilde f)\geq0\) on
\(\mathrm{PW}^{\mathrm{hol}}_F\).
\end{theorem}

\begin{proof}
The left-hand side has coefficients \(\Lambda(n)n^{-1/2}\) in the
additive variable \(\log n\). The Rankin--Selberg norm of a fixed \(F\)
has diagonal coefficients \(|a_F(n)|^2\) after unfolding. Equality for
all Paley--Wiener tests would force these coefficient systems to agree
as distributions, which they do not. Positivity of Petersson norm is
therefore irrelevant to the sign of the zeta finite-place Weil term
unless a separate trace formula is inserted.
\end{proof}

\begin{corollary}[Li approximants do not correct the prime kernel mismatch]
\label{v4-arith:w5-a:cor:HPLi-restricted}
\ClaimStatusProvedHere. Even if a Li test admits a
\(\mathrm{PW}^{\mathrm{hol}}_{F}\)-representative, the non-negativity of
\(W_{\mathrm{fin}}(\phi*\widetilde\phi)\) is not a consequence of the
fixed-\(F\) Rankin--Selberg norm. It requires an independent trace
formula for the von Mangoldt kernel.
\end{corollary}

\begin{proof}
This is exactly the obstruction isolated in
\Cref{v4-arith:w5-a:thm:prime-side}. The Li test changes the test
function; it does not change the kernel from \(\Lambda(n)\) to
\(|a_F(n)|^2\).
\end{proof}

\section{The Irreducible Obstruction: Archimedean Tate Distribution and Eisenstein}
\label{v4-arith:w5-a:sec:obstruction}

The fixed-\(F\) Rankin--Selberg norm does not identify
\(W_{\mathrm{fin}}(f*\widetilde{f})\) with a Petersson norm.
The archimedean Weil contribution \(W_{\infty}(f*\widetilde{f})\)
remains outside this identification: the Petersson archimedean
density is \(|\Gamma(\tfrac{1}{4}+\tfrac{it}{2})|^{2}\), while
\(W_{\infty}\) pairs against
\(\mathrm{Re}\,\psi(\tfrac{1}{4}+\tfrac{it}{2})=\partial_{s}\log|\Gamma(s/2+1/4)|^{2}|_{s=it}\),
the logarithmic derivative. The positivity of \(W(f*\widetilde{f})\)
is the full Weil inequality; the archimedean term, the finite
von Mangoldt term, and the two polar terms cannot be separated into
independent positivity assertions.

\subsection{Archimedean Tate Distribution}
\label{v4-arith:w5-a:obstruction:archimedean}

\begin{definition}[archimedean Tate distribution]
\label{v4-arith:w5-a:def:tate-archimedean}
\ClaimStatusDefinitional. The archimedean Weil contribution
\(W_{\infty}\) of \eqref{v4-arith:w5-a:eq:weil-archimedean} is the
pairing of \(\widehat{f}\) against the distribution
\begin{equation}
\label{v4-arith:w5-a:eq:tate-archimedean-distribution}
\mathcal{T}_{\infty}(t)\;:=\;\frac{1}{2\pi}\,\mathrm{Re}\,\psi\!\left(\tfrac{1}{4}+\tfrac{it}{2}\right)+\frac{\log\pi}{2\pi}.
\end{equation}
This is the archimedean \(\mathbb{R}\)-factor of the completed zeta
function: Tate's thesis (Tate 1950) assigns to each place \(v\) of
\(\mathbb{Q}\) a local zeta distribution, and the archimedean factor
\(\Lambda_{\infty}(s)=\pi^{-s/2}\Gamma(s/2)\) contributes
\(\mathcal{T}_{\infty}\) in the sign convention of
\eqref{v4-arith:w5-a:eq:weil-archimedean}.
\end{definition}

\begin{proposition}[archimedean distribution outside a fixed Rankin--Selberg norm]
\label{v4-arith:w5-a:prop:archimedean-not-absorbed}
\ClaimStatusProvedHere. The coefficientwise lift
\(\Upsilon_{F}\colon\mathrm{PW}^{\mathrm{hol}}_{F}\to\mathcal{V}^{\mathrm{prim}}_{\mathrm{cusp,temp}}\)
of \Cref{v4-arith:w5-a:def:upsilon} does not absorb the archimedean
distribution \(\mathcal{T}_{\infty}\). Specifically, the restriction of
the Petersson form to the archimedean factor
\(\mathcal{H}^{(\infty)}\subset\mathcal{V}^{\mathrm{prim}}\) pairs with
\(|\widehat{f}(t)|^{2}\) against the \emph{\(|\Gamma(\tfrac{1}{4}+\tfrac{it}{2})|^{2}\)}
density, while \(\mathcal{T}_{\infty}\) pairs against the
logarithmic-derivative density \(\mathrm{Re}\,\psi(\tfrac{1}{4}+\tfrac{it}{2})\);
the two densities are not proportional.
\end{proposition}

\begin{proof}
At \(\infty\), \(\mathcal{H}^{(\infty)}\) is the standard Heisenberg Fock
on \(\mathbb{C}[\alpha_{-1},\alpha_{-2},\ldots]|0\rangle_{\infty}\)
(Frenkel--Ben-Zvi \S3.5). The Petersson-restriction at \(\infty\) uses
the archimedean local Rankin--Selberg integral, which by
Jacquet 1972 SLN~260 \S4 (archimedean local factors for \(GL_{2}\))
evaluates to a Gamma-absolute-value kernel
\(|\Gamma(\tfrac{1}{4}+\tfrac{it}{2})|^{2}\cdot|\widehat{f}(t)|^{2}\).
This is positive and forms part of \(K_{F}(s)\) in
\eqref{v4-arith:w5-a:eq:KF-kernel}.

The Weil archimedean distribution \(\mathcal{T}_{\infty}\) pairs
\(|\widehat{f}(t)|^{2}\) against
\(\mathrm{Re}\,\psi(\tfrac{1}{4}+\tfrac{it}{2})\), the logarithmic
derivative of \(\Gamma(\tfrac{1}{4}+\tfrac{it}{2})\). The relation
between the Gamma absolute value and the digamma real part is
\(|\Gamma(\tfrac{1}{4}+\tfrac{it}{2})|^{2}\cdot\frac{d}{dt}\mathrm{arg}\,\Gamma(\tfrac{1}{4}+\tfrac{it}{2})=\mathrm{Im}\,\psi(\tfrac{1}{4}+\tfrac{it}{2})\cdot|\Gamma(\tfrac{1}{4}+\tfrac{it}{2})|^{2}\),
not the real part. Hence the Petersson archimedean density and
\(\mathcal{T}_{\infty}\) differ; the former is manifestly positive, the
latter is not sign-definite (it takes both signs as \(t\to\infty\) by
Stirling: \(\mathrm{Re}\,\psi(\tfrac{1}{4}+\tfrac{it}{2})\sim\frac{1}{2}\log(t^{2}/4)\)
for \(|t|\to\infty\), positive for large \(t\); but is negative near
\(t=0\)).
\end{proof}

\begin{remark}[archimedean Tate density and conformal anomaly]
\label{v4-arith:w5-a:rem:archimedean-obstruction}
The archimedean digamma contribution \(\mathcal{T}_{\infty}\) is the
Tate-local density for the real place: Tate's thesis assigns to
\(\Lambda_{\infty}(s)=\pi^{-s/2}\Gamma(s/2)\) the local factor whose
logarithmic derivative is
\(\frac12\psi(s/2)-\frac12\log\pi\); after moving the local term to
the side used in \eqref{v4-arith:w5-a:eq:weil-explicit}, the
distribution is written as in
\eqref{v4-arith:w5-a:eq:tate-archimedean-distribution}.
The identity \(\psi=\Gamma'/\Gamma=\partial_{s}\log\Gamma\) gives
\(\mathcal{T}_{\infty}(t)=\frac{1}{2\pi}\mathrm{Re}\,\psi(\tfrac14+\tfrac{it}{2})+\frac{\log\pi}{2\pi}\),
the logarithmic derivative of the archimedean Gamma factor evaluated
on the critical line. The Petersson weight \(y^{k-2}\) in
\eqref{v4-arith:w5-a:eq:petersson} pairs against the modular form
and produces \(|\Gamma(\tfrac{1}{4}+\tfrac{it}{2})|^{2}\), the
square absolute value, not its logarithmic derivative. The two
quantities are related by the Bohr--Mollerup identity, but they
differ as distributions on \(\mathbb{R}\); the archimedean
conformal-anomaly contribution at (\Cref{v4-arith:polyakov}) is the
logarithmic-derivative side, outside the Petersson integral's reach.
\end{remark}

\subsection{Eisenstein Continuum and Non-Tempered Locus}
\label{v4-arith:w5-a:obstruction:eisenstein}

\begin{proposition}[the native continuum comparison]
\label{v4-arith:w5-a:prop:eisenstein-residual}
The local Gibbs construction and its scalar ratios do not determine a
native Eisenstein pairing. The continuum comparison requires the native
maps and regularized spectral measures described in
\Cref{v4-arith:kms-gns:cor:eisenstein-final}. A Li-positivity argument
using a discrete automorphic core additionally needs a map or projection
to that core and an identity for the relevant Li trace.
\end{proposition}
\begin{proof}
The local Gibbs carrier is a Hilbert--Schmidt operator space. Its
partition ratio has no induced-representation argument. Thus it cannot
define the proposed native spectral density without an additional map.
The scalar normalization test of
\Cref{v4-arith:kms-gns:rem:match-is-FE} refutes one proposed identity;
it does not compute the actual native density. A Mellin transform alone
also supplies no projection to a native discrete automorphic subspace.
\end{proof}

\begin{remark}[non-tempered residual]
\label{v4-arith:w5-a:rem:non-tempered}
For Maass cusp forms, temperedness is conditional on the Ramanujan
conjecture for \(GL_{2}/\mathbb{Q}\); the current best bound is
\(7/64\) (Kim--Sarnak 2003). The non-tempered locus contributes a
controlled correction to \(K_{F}(s)\) at
\(\mathrm{Re}\,s=1/2\pm\theta\) for \(\theta\leq 7/64\), but the
Rankin--Selberg argument gives no positive-definite form there. The
positive automorphic Rankin--Selberg statement on non-tempered Maass
forms is therefore conditional on the full
Ramanujan conjecture.
This does not address the zeta finite-place Weil term.
\end{remark}

\section{Restricted-Class Theorem: HP--Li on the Holomorphic Cuspidal Tempered Class}
\label{v4-arith:w5-a:restricted-theorem}

\begin{theorem}[HP--Li restricted-class identity]
\label{v4-arith:w5-a:thm:HPLi-restricted}
\ClaimStatusProvedHere. Fix a normalised holomorphic cuspidal Hecke
eigenform \(F\in S_{k}(\mathrm{SL}_{2}(\mathbb{Z}))\) of weight
\(k\geq 12\). For any \(f\in\mathrm{PW}^{\mathrm{hol}}_{F}\) satisfying
(H1)--(H3) of \Cref{v4-arith:w5-a:lem:upsilon-cuspidal}, the
symmetrised test function \(f*\widetilde{f}\) satisfies the identity
\begin{equation}
\label{v4-arith:w5-a:eq:restricted-HPLi}
\widehat{f*\widetilde{f}}\bigl(\tfrac{1}{2i}\bigr)+\widehat{f*\widetilde{f}}\bigl(-\tfrac{1}{2i}\bigr)
-W_{\mathrm{fin}}(f*\widetilde{f})
\;=\;W(f*\widetilde f)+W_{\infty}(f*\widetilde{f}).
\end{equation}
Equivalently: the sum of polar contributions minus the finite-prime
term splits into the Weil functional and the archimedean Tate
distribution. This is only the Weil explicit formula rearranged; it is
not a positivity theorem.
\end{theorem}

\begin{proof}
The Weil explicit formula
\eqref{v4-arith:w5-a:eq:weil-explicit} applied to \(f*\widetilde{f}\)
gives
\[
W(f*\widetilde f)=
\widehat{f*\widetilde{f}}(\tfrac{1}{2i})
+\widehat{f*\widetilde{f}}(-\tfrac{1}{2i})
-W_{\infty}(f*\widetilde f)-W_{\mathrm{fin}}(f*\widetilde f).
\]
Moving \(W_{\infty}\) to the other side gives
\eqref{v4-arith:w5-a:eq:restricted-HPLi}.
\end{proof}

\begin{remark}[H3 is removable]
\label{v4-arith:w5-a:rem:H3-removable}
Hypothesis (H3) (double vanishing of the Mellin transform at \(s=0,1\))
is an artifact of the normalisation
\(\xi_{\mathbb{R}}(s)=\frac{1}{2}s(s-1)\Lambda_{\zeta}(s)\): working
directly with \(\xi_{\mathbb{R}}\) instead of \(\Lambda_{\zeta}\) absorbs
the Tate poles. Therefore (H3) is removable in the sense that its
violation is quantifiable: if \(\widetilde{f}_{M}(s)\) has zeros of
order \((a,b)\leq(2,2)\) at \(s=0,1\), the polar contribution to
\(W(f*\widetilde{f})\) is corrected by a \(2(2-a)+2(2-b)\)-fold Taylor
expansion at \(s=0,1\), each term a specific derivative of
\(\xi_{\mathbb{R}}\) at these points. The correction is a finite-rank
bounded operator on \(\mathrm{PW}(\mathbb{R})\), hence does not affect
the positivity question on the main spectral locus.
\end{remark}

\begin{remark}[H2 is the genuine obstruction]
\label{v4-arith:w5-a:rem:H2-obstruction}
Hypothesis (H2) (holomorphy of \(\widetilde{f}_{M}\) on the strip of
width \((k-1)\)) is the genuine restriction. In the limit
\(k\to\infty\) the strip widens to all of \(\mathbb{C}\), but no single
finite-weight \(F\) gives access to the full Paley--Wiener cone. The
full cone requires either (i) a non-holomorphic lift to Maass forms
(where \(K_{F}\)'s positivity becomes conditional on Ramanujan), or
(ii) a limiting procedure as \(k\to\infty\) which is problematic at
the level of the Petersson measure (the Plancherel weight vanishes).
The limiting procedure is precisely the Selberg trace formula in
disguise, and its positivity is Weil positivity itself.
\end{remark}

\section{The Archimedean Tate Positivity Conjecture}
\label{v4-arith:w5-a:irreducible-obstruction}

\begin{conjecture}[archimedean Tate positivity, A-TP]
\label{v4-arith:w5-a:conj:archimedean-tate-positivity}
For every \(f\in\mathrm{PW}(\mathbb{R})\),
\begin{equation}
\label{v4-arith:w5-a:eq:archimedean-tate-positivity}
W_{\infty}(f*\widetilde{f})\;\leq\;
\widehat{f*\widetilde{f}}(\tfrac{1}{2i})+
\widehat{f*\widetilde{f}}(-\tfrac{1}{2i})-
W_{\mathrm{fin}}(f*\widetilde{f}).
\end{equation}
Equivalently, \(W(f*\widetilde{f})\geq0\). Equivalently: RH.
\end{conjecture}

\begin{theorem}[reformulation of A-TP]
\label{v4-arith:w5-a:thm:ATP-equals-weil-positivity}
\ClaimStatusProvedHere. Conjecture A-TP
(\Cref{v4-arith:w5-a:conj:archimedean-tate-positivity}) is strictly
equivalent to Weil positivity
\Cref{v4-arith:w5-a:cor:weil-positivity-RH}, and hence to RH.
\end{theorem}

\begin{proof}
By \eqref{v4-arith:w5-a:eq:weil-explicit},
\[
\widehat{f*\widetilde f}(\tfrac{1}{2i})+
\widehat{f*\widetilde f}(-\tfrac{1}{2i})-
W_{\mathrm{fin}}(f*\widetilde f)-
W_{\infty}(f*\widetilde f)
=W(f*\widetilde f).
\]
Thus \eqref{v4-arith:w5-a:eq:archimedean-tate-positivity} is exactly
\(W(f*\widetilde f)\geq0\). The equivalence to RH is
\Cref{v4-arith:w5-a:cor:weil-positivity-RH}.
\end{proof}

\begin{remark}[domain of the Mellin--Rankin--Selberg isometry]
\label{v4-arith:w5-a:rem:honest-form}
The Mellin--Rankin--Selberg isometry gives positivity for automorphic
Rankin--Selberg kernels. It does not determine the zeta finite-place Weil
distribution. The residual is therefore not merely the sign of
\(W_{\infty}\); the full residual is the inequality
\[
W_{\infty}(f*\widetilde f)\le
\widehat{f*\widetilde f}(\tfrac{1}{2i})+
\widehat{f*\widetilde f}(-\tfrac{1}{2i})-
W_{\mathrm{fin}}(f*\widetilde f),
\]
which is exactly Weil positivity.
\end{remark}

\section{HP--Li Boundary}
\label{v4-arith:w5-a:form}

The decomposition of HP--Li positivity separates identities from
positivity. The prime-side Weil contribution
\(W_{\mathrm{fin}}(f*\widetilde{f})\) is not made positive by the fixed
Rankin--Selberg lift of \Cref{v4-arith:w5-a:thm:prime-side}. The polar
terms \(\widehat{f*\widetilde{f}}(\pm 1/(2i))\) are explicit finite
Taylor corrections (\Cref{v4-arith:w5-a:rem:H3-removable}).
The full A-TP inequality with \(W_{\infty}\) isolated is RH-equivalent
(\Cref{v4-arith:w5-a:conj:archimedean-tate-positivity},
\Cref{v4-arith:w5-a:thm:ATP-equals-weil-positivity}). The extension to
Maass cuspidal tempered forms is conditional on the Ramanujan conjecture
(\Cref{v4-arith:w5-a:rem:non-tempered}). The Eisenstein continuous
spectrum is domain-excluded by the density mismatch of
\Cref{v4-arith:w5-a:prop:eisenstein-residual}. The full
\(\mathrm{PW}(\mathbb{R})\) class is obstructed by (H2)
(\Cref{v4-arith:w5-a:rem:H2-obstruction}).

\begin{theorem}[HP--Li positivity boundary after the kernel check]
\label{v4-arith:w5-a:thm:hpli-form}
\ClaimStatusProvedHere. HP--Li positivity
(\Cref{v4-arith:rh:thm:VB-implies-F-RH}) does not follow from the
fixed-\(F\) Rankin--Selberg lift. That lift gives an automorphic
comparison, but it does not identify the von Mangoldt finite-place
kernel with a Petersson norm. The RH-equivalent positivity statement is
the full archimedean Tate inequality
\Cref{v4-arith:w5-a:conj:archimedean-tate-positivity}, equivalently
Weil positivity, equivalently RH.
\end{theorem}

\begin{proof}
Assemble:
\Cref{v4-arith:w5-a:thm:prime-side},
\Cref{v4-arith:w5-a:thm:RS-integral},
\Cref{v4-arith:w5-a:thm:HPLi-restricted},
\Cref{v4-arith:w5-a:prop:archimedean-not-absorbed},
\Cref{v4-arith:w5-a:thm:ATP-equals-weil-positivity}.
The first two results identify the kernel mismatch; the third is only
an algebraic rearrangement of Weil's formula; the final result records
the RH-equivalent positivity statement.
\end{proof}

\begin{corollary}[HP--Li inputs in the sufficient RH path]
\label{v4-arith:w5-a:cor:sufficient-path-impact}
\ClaimStatusProvedHere. The sufficient RH path of
\Cref{v4-arith:rh:thm:sufficient-path} requires the following HP--Li
inputs:
\begin{itemize}
\item a trace formula for the zeta von Mangoldt kernel beyond the
fixed-\(F\) lift;
\item the A-TP inequality, which is RH-equivalent
(\Cref{v4-arith:w5-a:conj:archimedean-tate-positivity}).
\item an extension beyond \(\mathrm{PW}^{\mathrm{hol}}_{F}\), either
through Maass lifts with Ramanujan control or through a
weight-\(k\to\infty\) limit with the Petersson limit measure under
control.
\end{itemize}
The RH-equivalent residual among these inputs is A-TP:
\[
W_{\infty}(f*\widetilde f)\le
\widehat{f*\widetilde f}(\tfrac{1}{2i})+
\widehat{f*\widetilde f}(-\tfrac{1}{2i})-
W_{\mathrm{fin}}(f*\widetilde f),
\]
which is equivalent to RH.
The prime-side comparison and the extension beyond
\(\mathrm{PW}^{\mathrm{hol}}_{F}\) remain separate domain conditions.
The HP--Li obstruction is the full Weil inequality with finite, polar,
and archimedean terms kept in their correct places.
\end{corollary}

\begin{proof}
Combine \Cref{v4-arith:w5-a:thm:hpli-form} with the sufficient-path
theorem \Cref{v4-arith:rh:thm:sufficient-path}.
\end{proof}

\begin{remark}[decomposition of HP--Li positivity]
\label{v4-arith:w5-a:rem:domain-summary}
The kernel identity separates Weil positivity into four constituents:
\begin{enumerate}
\item The finite-prime term
\(W_{\mathrm{fin}}(f*\widetilde{f})\): a von Mangoldt distribution not
identified here with a fixed-\(F\) Petersson norm.
\item The archimedean term: the Tate
\(\mathcal{T}_{\infty}\)-distribution
\(\psi(\tfrac{1}{4}+\tfrac{it}{2})=\partial_{s}\log\Gamma|_{s=1/4+it/2}\),
which is sign-indefinite.
\item The decomposition identity
\(W(f*\widetilde{f})=\widehat{f*\widetilde f}(\tfrac{1}{2i})
+\widehat{f*\widetilde f}(-\tfrac{1}{2i})
-W_{\mathrm{fin}}(f*\widetilde f)-W_{\infty}(f*\widetilde f)\).
\item The test-function class \(\mathrm{PW}^{\mathrm{hol}}_{F}\),
a restricted subspace of \(\mathrm{PW}(\mathbb{R})\), on which the
full HP--Li inequality is not implied by the fixed-\(F\) lift.
\end{enumerate}
HP--Li on the full Paley--Wiener class is equivalent to RH.
\end{remark}

\section{Beilinson Boundary}
\label{v4-arith:w5-a:summary}

\begin{enumerate}
\item The coefficientwise Mellin lift
\(\Upsilon_{F}\colon\mathrm{PW}^{\mathrm{hol}}_{F}\to q\mathbb C[[q]]\)
of \Cref{v4-arith:w5-a:def:upsilon} is a formal probe; it is not a
modular-form-valued map for an arbitrary Paley--Wiener multiplier.
\item The Rankin--Selberg kernel
\Cref{v4-arith:w5-a:thm:RS-integral} is built from
\(L(s,F\otimes\overline F)\), not from \(-\zeta'/\zeta\).
\item The prime-side Weil contribution is not proved positive by the
fixed-\(F\) Petersson norm
(\Cref{v4-arith:w5-a:thm:prime-side}).
\item The irreducible obstruction to full HP--Li is the full A-TP
inequality, including polar, finite, and archimedean terms; it is
strictly equivalent to Weil positivity and to RH
(\Cref{v4-arith:w5-a:conj:archimedean-tate-positivity},
\Cref{v4-arith:w5-a:thm:ATP-equals-weil-positivity}).
\item The Eisenstein continuous spectrum is domain-excluded by the
density mismatch of \Cref{v4-arith:kms-gns:cor:eisenstein-final}; it
does not contribute to the Li coefficients, which see only the
discrete spectrum.
\item HP--Li form
(\Cref{v4-arith:w5-a:thm:hpli-form}): the full HP--Li inequality is
RH-equivalent.
\end{enumerate}

\begin{remark}[Beilinson principle: restricted-class HP--Li]
\label{v4-arith:w5-a:rem:beilinson-principle}
The restricted-class HP--Li theorem is an identity together with a
kernel obstruction. The full claim, HP--Li on all of
\(\mathrm{PW}(\mathbb{R})\), is equivalent to RH and remains open as
Conjecture A-TP (\Cref{v4-arith:w5-a:conj:archimedean-tate-positivity}).
The lift \(\Upsilon_F\) is a useful formal comparison, not a proof that
the von Mangoldt kernel is a positive Petersson norm.
\end{remark}

\section{Bibliography of Primary Sources}
\label{v4-arith:w5-a:bibliography}

The identities above use the following primary sources, disjoint from
the proof sources used in
Chapters~\ref{v4-arith:closure} (P1--P3 resolution),
\ref{v4-arith:kp-compat} (Koszul--Pet compatibility), and
\ref{v4-arith:kms-gns-full} (KMS\({}_{1}\)/GNS full sector).

\begin{description}
\item[Weil 1952.] Weil, A. \emph{Sur les formules explicites de la
th\'eorie des nombres premiers}. \emph{Meddelanden fr\r{a}n Lunds
Univ.~Mat.~Sem.} (Festschrift f\"or Marcel Riesz), 1952; reprinted
\emph{Oeuvres Scientifiques} Vol.~II, Springer 1979. Explicit formula,
\(W\)-distribution, Weil positivity equivalence to RH.
\item[Li 1997.] Li, X.-J. \emph{The positivity of a sequence of
numbers and the Riemann hypothesis}. \emph{J.~Number Theory} 65
(1997), 325--333. Li coefficients \(\lambda_{n}\),
equivalence to RH, explicit Paley--Wiener realisation.
\item[Rankin 1939.] Rankin, R.~A. \emph{Contributions to the theory
of Ramanujan's function \(\tau(n)\) and similar arithmetical functions}.
\emph{Proc.~Camb.~Phil.~Soc.}~35 (1939), 357--372. Rankin--Selberg
identity for \(L\)-series of modular forms.
\item[Selberg 1940.] Selberg, A. \emph{Bemerkungen \"uber eine
Dirichletsche Reihe, die mit der Theorie der Modulformen nahe
verbunden ist}. \emph{Arch.~Math.~Naturvid.}~43 (1940), 47--50.
Rankin--Selberg identity and \(L\)-function analytic continuation.
\item[Hecke 1936.] Hecke, E. \emph{\"Uber die Bestimmung Dirichletscher
Reihen durch ihre Funktionalgleichung}. \emph{Math.~Ann.} 112
(1936), 664--699. Mellin transform / modular form correspondence
and functional equation.
\item[Deligne 1974.] Deligne, P. \emph{La conjecture de Weil I}.
\emph{Publ.~Math.~IHES} 43 (1974), 273--307. Ramanujan bound for
holomorphic cuspidal Hecke eigenforms of weight \(k\geq 2\).
\item[Jacquet 1972.] Jacquet, H. \emph{Automorphic forms on \(GL_{2}\)
II}. Lecture Notes in Mathematics 278, Springer 1972. Local
archimedean factors for \(GL_{2}\); ramified principal series.
\item[Jacquet--Langlands SLN~114 1970.] Jacquet, H.;
Langlands, R.~P. \emph{Automorphic forms on \(GL(2)\)}. Lecture Notes
in Mathematics 114, Springer 1970. Adelic Whittaker model, Hecke
theory.
\item[Iwaniec--Kowalski 2004.] Iwaniec, H.; Kowalski, E.
\emph{Analytic Number Theory}. AMS Colloquium Publications 53, 2004.
Modern account of Weil's explicit formula, Rankin--Selberg,
Paley--Wiener.
\item[Bombieri--Lagarias 1999.] Bombieri, E.; Lagarias, J.~C.
\emph{Complements to Li's criterion for the Riemann hypothesis}.
\emph{J.~Number Theory} 77 (1999), 274--287. Alternative constructions
of Li-coefficient approximants and equivalences to Weil positivity.
\item[Connes 1999.] Connes, A. \emph{Trace formula in noncommutative
geometry and the zeros of the Riemann zeta function}. \emph{Selecta
Math.~(N.~S.)} 5 (1999), 29--106, arXiv:math/9811068. Adele class
space, absorption trace, Weil positivity as trace-formula positivity.
\item[Patterson 1988.] Patterson, S.~J. \emph{An Introduction to the
Theory of the Riemann Zeta-Function}. Cambridge Studies in Advanced
Mathematics 14, 1988. Complete treatment of the Hadamard product
and explicit formula.
\item[Kim--Sarnak 2003.] Kim, H.~H.; Sarnak, P.~C. Appendix to
Kim, H.~H., \emph{Functoriality for the exterior square of \(GL_{4}\)
and the symmetric fourth of \(GL_{2}\)}. \emph{J.~Amer.~Math.~Soc.} 16
(2003), 139--183. Bound \(|a_{F}(p)|\leq 2p^{7/64}\) towards the
Ramanujan conjecture for \(GL_{2}/\mathbb{Q}\).
\item[Bump 1997.] Bump, D. \emph{Automorphic Forms and
Representations}. Cambridge Studies in Advanced Mathematics 55, 1997.
Comprehensive account of Rankin--Selberg, Whittaker models,
Petersson theory.
\item[Zagier 1981.] Zagier, D. \emph{The Rankin--Selberg method for
automorphic functions which are not of rapid decay}.
\emph{Bonner Math.~Schriften} 121 (1981). Rankin--Selberg unfolding
technique.
\end{description}

\noindent\emph{Disjointness comparison.} The primary-source list
for the identities above (Weil 1952 for the explicit formula; Li 1997 for
Li coefficients; Rankin 1939 and Selberg 1940 for the Rankin--Selberg
identity; Hecke 1936 for Mellin--functional-equation duality;
Deligne 1974 for temperedness; Jacquet 1972 and Jacquet--Langlands
SLN~114 for adelic \(GL_{2}\) representation theory) is disjoint from
the lists underwriting P1 (Beilinson--Drinfeld 2004, Positselski
2009), P2 (Tate 1950), P3 / Koszul--Pet (Bost--Connes 1995,
Takesaki 1979, Jacquet--Piatetski-Shapiro--Shalika 1983,
Godement--Jacquet SLN~260, Moeglin--Waldspurger 1995, Arthur 1980,
Shahidi 2010), the Tate restricted tensor chapter (Lurie HA,
Fresse~2017, Ayala--Francis), and the HP\(^{\mathrm{Ar}}\) / VB\(^{*}\)
converse lists (Connes 1999 is a cross-reference cited for
context in both arguments; it does not appear on both sides of any
single identification).
